\section{Introduction}

Connected laboratories are assembled from instruments, vendor software, open-source libraries, physical-resource definitions, schedulers, sample records, and local operating practice. Their maintenance depends on work that formal system descriptions rarely capture. A practitioner must identify which interface is actually available, determine whether a file is portable source or active runtime state, reconstruct which physical actions completed before an exception, establish whether a telemetry pattern indicates failure, and negotiate what evidence would count as resolution. This work crosses organizational and representational boundaries. It is social, technical, and materially consequential.

Public specialist forums provide one setting in which this maintenance becomes inspectable. A user posts an error or design question. Respondents request versions, layouts, liquid properties, licence status, traces, sample constraints, and prior changes. Vendor representatives clarify product behavior. Open-source maintainers translate incidents into issues or corrections. Participants sometimes move into direct messages, private support, meetings, or local implementations, creating a boundary between visible discussion and hidden resolution. In stronger cases, the community converts a local incident into a public post-mortem, example, schema, or maintained code change.

CSCW research has long examined the articulation work required to coordinate distributed activity, the infrastructural conditions that become visible during breakdown, and repair as a process that sustains sociotechnical order \cite{schmidt1992articulation,star1996infrastructure,jackson2014repair}. Laboratory automation provides a demanding case. Coordination must align software state with instruments, material location, sample value, calibration, timing, contamination history, and scientific acceptance. A software exception can occur after an irreversible physical action. A public answer can be technically correct and still fail to become durable community knowledge. An interface can exist and remain operationally inaccessible.

This paper examines a purposive corpus of public laboratory-automation discussions to ask:
\begin{enumerate}
  \item How do participants articulate hidden technical and scientific dependencies during public troubleshooting and design work?
  \item How is repair coordinated when software histories, physical actions, and material state diverge?
  \item Under what conditions does local problem solving become durable public infrastructure, and where does it migrate beyond the visible community record?
\end{enumerate}

The analysis contains 55 selected discussions and 45 episodes from a deliberately difficult 14-thread subset. A ten-construct coding model provides a structured map of knowledge, interfaces, representations, runtime coordination, validation, and AI context. The CSCW contribution lies in the relational structure across those constructs. We show that public technical support operates as distributed articulation work, that recovery after partial execution is a form of material repair requiring evidence about an irreversible prefix, and that knowledge durability depends on infrastructuring practices that attach provenance, applicability, ownership, validation, and correction pathways to an answer.

The paper contributes three elements. First, it offers a boundary-centered account of maintenance in connected laboratories. Second, it identifies public/private migration, validation-stage negotiation, and final material disposition as central features of collaborative repair in cyber-physical work. Third, it derives design implications for question interfaces, support systems, event records, and community-maintained knowledge infrastructures. The study does not treat discussion frequency as operational incidence and does not support vendor ranking or causal inference.
